%% ch07.tex — Chapter 7: SQL at Scale: HiveQL and the SQL-on-Hadoop Ecosystem
%% Part of "Data Engineering" textbook

\chapter{SQL at Scale: HiveQL and the SQL-on-Hadoop Ecosystem}
\label{ch:sql-scale}

\epigraph{\textit{``Presto is the democratisation of data.
Before Presto, only engineers could query our data warehouse.
After Presto, every product manager, every analyst, every
curious Facebooker could answer their own questions.''}}{Dain Sundstrom, Presto co-creator, Meta Engineering Blog, 2013}

\vspace{4pt}

\begin{chapterlearning}
After completing this chapter, you will be able to:
\begin{enumerate}
  \item Use HiveQL complex data types---\texttt{ARRAY}, \texttt{MAP},
        and \texttt{STRUCT}---to model semi-structured data, and apply
        \texttt{LATERAL VIEW EXPLODE} to unnest array and map columns
        into relational rows.
  \item Write SQL window function queries using \texttt{OVER},
        \texttt{PARTITION BY}, \texttt{ORDER BY}, \texttt{ROWS BETWEEN},
        and the ranking functions \texttt{ROW\_NUMBER}, \texttt{RANK},
        \texttt{LAG}, and \texttt{LEAD}.
  \item Distinguish the three Hive UDF categories---scalar UDF, User-Defined
        Aggregate Function (UDAF), and User-Defined Table-Generating Function
        (UDTF)---and identify the appropriate category for a given extension.
  \item Explain how Apache Tez improves on Hive-on-MapReduce by replacing
        the chain of MapReduce jobs with a single DAG of vertices and
        edges, and identify which types of intermediate materialisation
        Tez eliminates.
  \item Describe the Hive LLAP architecture---the persistent daemon,
        in-memory columnar cache, and shared executor threads---and explain
        why LLAP enables sub-second interactive query latency.
  \item Explain how the Cost-Based Optimizer (CBO) uses table and column
        statistics from the Hive Metastore to choose join order and join
        strategy, and describe the difference between rule-based and
        cost-based optimisation.
  \item Compare Hive's four join execution strategies---Common Join,
        Map Join, Bucket Map Join, and Sort-Merge Bucket Join---and select
        the correct strategy for a given table size and bucketing configuration.
  \item Describe the Presto/Trino architecture---Coordinator, Workers,
        and Connectors---and explain how the connector model enables
        federated SQL queries across heterogeneous data sources.
  \item Explain the key design differences between Hive and Presto that
        make Presto suitable for interactive queries ($<$1 minute) and
        Hive suitable for long-running batch ETL (hours).
  \item Compare the four major SQL-on-Hadoop engines---Hive, Presto/Trino,
        Spark SQL, and Apache Impala---on latency, scale, fault tolerance,
        and data source support, and select the appropriate engine for
        a given workload.
\end{enumerate}
\end{chapterlearning}

\bigskip

Chapter~\ref{ch:mapreduce} introduced Apache Hive as the SQL layer that
made Hadoop accessible to analysts without MapReduce expertise, covering
the Hive Metastore, partitioning, bucketing, and the basic HiveQL-to-MapReduce
compilation pipeline. Chapter~\ref{ch:spark} showed how Spark SQL and the
Catalyst optimizer brought a faster, in-memory SQL engine to the same Hive
Metastore and Parquet ecosystem.

This chapter completes the SQL-at-scale picture by examining what sits
between those two endpoints: the advanced features of HiveQL that make
it a production-grade analytical SQL dialect; the successive execution
engines (Tez, LLAP) that have reduced Hive query latency from hours to
seconds; the cost-based optimizer that brings intelligent query planning
to Hive; and the family of alternative SQL engines---most importantly
Presto/Trino---that address the interactive analytics use case that neither
Hive-on-MapReduce nor batch Spark optimally serves.

The central theme is \emph{trade-off}: no single SQL engine is optimal
for all workloads. Understanding what each engine optimises for---and
what it sacrifices in return---is the core competency this chapter develops.

%%====================================================================
\section{HiveQL: Advanced Language Features}
\label{sec:hiveql-advanced}
%%====================================================================

HiveQL extends ANSI SQL with features designed specifically for
semi-structured data and large-scale analytical workloads. Three categories
of extensions---complex data types, window functions, and user-defined
functions---are essential for production data engineering work.

%--------------------------------------------------------------------
\subsection{Complex Data Types: ARRAY, MAP, and STRUCT}
\label{subsec:complex-types}
%--------------------------------------------------------------------

Real-world data rarely arrives in fully normalised form. A web event
may carry a variable-length list of product categories, an e-commerce
order may embed a nested shipping address, and a user profile may contain
a key-value map of arbitrary feature flags. Hive models these patterns
with three native complex types.

An \term{ARRAY} is an ordered sequence of values of the same type. An
\term{MAP} is an unordered collection of key-value pairs. A \term{STRUCT}
is a named tuple of fields with potentially different types, analogous
to a JSON object:

\begin{lstlisting}[style=hiveql, caption={HiveQL: table schema with complex data types}]
CREATE TABLE user_events (
    event_id    BIGINT,
    user_id     BIGINT,
    event_ts    TIMESTAMP,
    -- ARRAY: a user may view 0..N product categories
    categories  ARRAY<STRING>,
    -- MAP: arbitrary key-value feature flags
    features    MAP<STRING, DOUBLE>,
    -- STRUCT: nested shipping address
    ship_addr   STRUCT<
                    street: STRING,
                    city:   STRING,
                    postal: STRING,
                    country: STRING
                >
)
STORED AS ORC
PARTITIONED BY (dt STRING);

-- Accessing nested fields
SELECT  event_id,
        categories[0]           AS first_category,
        features['ctr_score']   AS ctr_score,
        ship_addr.city          AS city,
        ship_addr.country       AS country
FROM    user_events
WHERE   dt = '2024-03-01'
  AND   size(categories) > 0;
\end{lstlisting}

%--------------------------------------------------------------------
\subsection{LATERAL VIEW and EXPLODE}
\label{subsec:lateral-view}
%--------------------------------------------------------------------

Aggregating or filtering over the elements of an ARRAY or MAP column
requires first \emph{exploding} the collection into separate rows.
\term{LATERAL VIEW EXPLODE} generates one output row per element of the
collection, joining it back to the containing row's other columns:

\begin{lstlisting}[style=hiveql, caption={HiveQL: LATERAL VIEW EXPLODE to unnest an array}]
-- Count page views by category across all events
SELECT  cat,
        COUNT(*)        AS views,
        COUNT(DISTINCT user_id) AS unique_users
FROM    user_events
LATERAL VIEW EXPLODE(categories) cat_table AS cat
WHERE   dt BETWEEN '2024-03-01' AND '2024-03-31'
GROUP BY cat
ORDER BY views DESC
LIMIT 20;
\end{lstlisting}

Without \texttt{LATERAL VIEW EXPLODE}, this query would require a
user-defined function to iterate over the array. With it, the HiveQL
compiler expands each \texttt{user\_events} row into as many rows as
there are elements in its \texttt{categories} array, making the standard
\texttt{GROUP BY} aggregation applicable directly.

\begin{defbox}{LATERAL VIEW EXPLODE}
\texttt{LATERAL VIEW EXPLODE(col)} generates a virtual table by
unnesting an \texttt{ARRAY} or \texttt{MAP} column. For an
\texttt{ARRAY<T>} column, each element produces one row. For a
\texttt{MAP<K,V>} column, each key-value pair produces one row
with two virtual columns (key and value). The virtual table is
cross-joined to the source row via the lateral view construct.
\end{defbox}

For \texttt{MAP} columns, \texttt{EXPLODE} produces a two-column
virtual table:

\begin{lstlisting}[style=hiveql, caption={HiveQL: EXPLODE on a MAP column}]
-- Find the average feature value by feature name
SELECT  feature_name,
        AVG(feature_val) AS avg_val
FROM    user_events
LATERAL VIEW EXPLODE(features) feat_table AS feature_name, feature_val
WHERE   dt = '2024-03-01'
GROUP BY feature_name;
\end{lstlisting}

%--------------------------------------------------------------------
\subsection{Window Functions}
\label{subsec:window-functions}
%--------------------------------------------------------------------

\term{Window functions} compute aggregate values over a sliding window
of rows around each output row, without collapsing the result set the
way \texttt{GROUP BY} does. They are declared with the \texttt{OVER}
clause and are essential for time-series analysis, ranking, sessionisation,
and cohort analytics.

The general syntax is:
\[
\texttt{function\_name(args) OVER ([PARTITION BY cols] [ORDER BY cols] [frame])}
\]

\begin{lstlisting}[style=hiveql, caption={HiveQL: window functions for sales ranking and trends}]
SELECT
    region,
    product_id,
    sale_date,
    revenue,

    -- Running total within region, ordered by date
    SUM(revenue) OVER (
        PARTITION BY region
        ORDER BY sale_date
        ROWS BETWEEN UNBOUNDED PRECEDING AND CURRENT ROW
    )   AS running_total,

    -- 7-day moving average
    AVG(revenue) OVER (
        PARTITION BY region, product_id
        ORDER BY sale_date
        ROWS BETWEEN 6 PRECEDING AND CURRENT ROW
    )   AS ma7,

    -- Rank product by revenue within region (ties get same rank)
    RANK() OVER (
        PARTITION BY region
        ORDER BY revenue DESC
    )   AS revenue_rank,

    -- Row number (no ties; deterministic ordering)
    ROW_NUMBER() OVER (
        PARTITION BY region
        ORDER BY revenue DESC
    )   AS row_num,

    -- Revenue of the previous day for the same product-region
    LAG(revenue, 1, 0) OVER (
        PARTITION BY region, product_id
        ORDER BY sale_date
    )   AS prev_day_revenue,

    -- Revenue of the next day
    LEAD(revenue, 1, 0) OVER (
        PARTITION BY region, product_id
        ORDER BY sale_date
    )   AS next_day_revenue

FROM daily_sales
WHERE dt = '2024-Q1';
\end{lstlisting}

\begin{table}[ht]
\centering
\caption{Key HiveQL window functions}
\label{tab:window-functions}
\renewcommand{\arraystretch}{1.5}
\begin{tabular}{L{3.2cm} L{7.2cm}}
\toprule
\textbf{Function} & \textbf{Description} \\
\midrule
\texttt{ROW\_NUMBER()}    & Unique sequential integer per row within partition; no ties \\
\texttt{RANK()}           & Rank with gaps for ties (e.g., 1, 2, 2, 4) \\
\texttt{DENSE\_RANK()}    & Rank without gaps (e.g., 1, 2, 2, 3) \\
\texttt{NTILE(n)}         & Divides partition into $n$ approximately equal buckets \\
\texttt{PERCENT\_RANK()}  & $(rank - 1) / (rows - 1)$; range [0, 1] \\
\texttt{LAG(col, n, def)} & Value of \texttt{col} from $n$ rows before the current row \\
\texttt{LEAD(col, n, def)} & Value of \texttt{col} from $n$ rows after the current row \\
\texttt{FIRST\_VALUE(col)} & First value in the window frame \\
\texttt{LAST\_VALUE(col)}  & Last value in the window frame \\
\texttt{SUM/AVG/MIN/MAX}  & Standard aggregates over the window frame \\
\bottomrule
\end{tabular}
\end{table}

%--------------------------------------------------------------------
\subsection{User-Defined Functions}
\label{subsec:udfs}
%--------------------------------------------------------------------

When built-in HiveQL functions are insufficient, analysts can extend
HiveQL through three categories of \term{User-Defined Functions} (UDFs):

\textbf{Scalar UDF:} maps one input row to one output row (one-to-one).
The most common category: URL parsing, custom hashing, text normalisation,
geolocation lookups. Implemented by extending Hive's
\texttt{org.apache.hadoop.hive.ql.exec.UDF} class in Java, or defined
inline as a temporary function from a Python/R script via HiveQL's
\texttt{TRANSFORM ... USING} clause.

\textbf{User-Defined Aggregate Function (UDAF):} maps many input rows to
one output row per group (many-to-one). Use when the built-in aggregates
(\texttt{SUM}, \texttt{AVG}, \texttt{COUNT}, \texttt{COLLECT\_LIST}) are
insufficient---for example, computing a custom weighted median or a sketch
(HyperLogLog approximate distinct count). UDAFs must implement both a
\texttt{partial} phase (within the Map task, analogous to a Combiner)
and a \texttt{final} phase (within the Reduce task).

\textbf{User-Defined Table-Generating Function (UDTF):} maps one input
row to zero or more output rows (one-to-many). This is the category that
\texttt{EXPLODE} belongs to: it takes one row with an array column and
produces one row per array element. Custom UDTFs are appropriate for
parsing XML or JSON strings embedded in a column, splitting a delimited
string into multiple rows, or expanding a compressed binary event log.

%%====================================================================
\section{HiveServer2 and Concurrent Access}
\label{sec:hiveserver2}
%%====================================================================

Original Hive (version 0.x) ran as a command-line tool that executed
one query at a time. For a production data warehouse serving hundreds of
concurrent analysts, this is untenable. \term{HiveServer2}, introduced
in Hive 0.11 (2013), is a long-running service that exposes Hive's query
capabilities over JDBC and ODBC interfaces, supporting concurrent sessions
from multiple clients simultaneously.

HiveServer2's architecture separates query compilation from execution.
A pool of \term{HiveServer2 daemon threads} each manage one client session,
compiling the client's HiveQL into a query plan and submitting it to the
cluster's execution engine. Multiple threads may submit queries to YARN
simultaneously, with the Fair Scheduler or Capacity Scheduler arbitrating
cluster resources between them.

\begin{systemdesign}{HiveServer2 in Production: Beeline and JDBC}
The standard command-line client for HiveServer2 is \texttt{beeline},
which connects over JDBC:
\begin{lstlisting}[style=shell]
beeline -u "jdbc:hive2://hs2-host:10000/analytics_db" \
        -n analyst_user --password=... \
        -e "SELECT COUNT(*) FROM web_events WHERE dt='2024-03-01'"
\end{lstlisting}
Business intelligence tools (Tableau, Power BI, Apache Superset) connect
to HiveServer2 via their Hive JDBC/ODBC drivers, making the full Hive
warehouse accessible through standard BI tooling without any
Hadoop-specific configuration on the analyst's workstation. HiveServer2
also supports Kerberos authentication for secure enterprise deployments.
\end{systemdesign}

%%====================================================================
\section{Execution Engine Evolution: MapReduce, Tez, and LLAP}
\label{sec:execution-engines}
%%====================================================================

The HiveQL language has remained stable since Hive 1.0, but the execution
engines that compile HiveQL into distributed computation have undergone
three successive architectural generations, each eliminating a specific
performance bottleneck of its predecessor.

%--------------------------------------------------------------------
\subsection{Hive-on-MapReduce: The Baseline}
\label{subsec:hive-mr}
%--------------------------------------------------------------------

Chapter~\ref{ch:mapreduce} established that a HiveQL query compiles to
a sequence of MapReduce jobs. A join followed by an aggregation compiles
to two jobs: Job~1 performs the join (Map reads input tables, Reduce
co-locates matching rows); Job~2 performs the aggregation (Map reads the
join output, Reduce computes the aggregate). Between Job~1 and Job~2,
the entire join result is materialised to HDFS.

This HDFS materialisation is the core performance problem. For a pipeline
with $k$ intermediate stages, each transition writes and reads the full
intermediate result from HDFS: $2 \times (k-1)$ HDFS operations. For a
complex ETL query with 5 stages, that is 8 full HDFS read/write cycles
before the final result is produced. At 1\TB\ of intermediate data per
stage, this amounts to 8\TB\ of unnecessary disk I/O in addition to the
actual computation.

%--------------------------------------------------------------------
\subsection{Apache Tez: DAG-Based Execution}
\label{subsec:tez}
%--------------------------------------------------------------------

\term{Apache Tez}, released in 2014 and developed primarily by Hortonworks,
replaces the MapReduce execution model with a \term{directed acyclic graph
(DAG)} of computation vertices connected by data edges. Tez allows Hive
to express a complex multi-stage query as a single DAG rather than a chain
of separate MapReduce jobs.

\begin{figure}[ht]
\centering
\begin{tikzpicture}[
    mr/.style={rectangle, draw=DERed!60, fill=DERed!15, rounded corners=4pt,
               minimum width=2.0cm, minimum height=0.65cm,
               font=\small\sffamily, align=center},
    tez/.style={rectangle, draw=DEGreen!70, fill=DELightGreen, rounded corners=4pt,
                minimum width=2.0cm, minimum height=0.65cm,
                font=\small\sffamily, align=center},
    disk/.style={cylinder, draw=DEGray!60, fill=DEGray!15, shape border rotate=90,
                 minimum height=0.55cm, minimum width=0.9cm,
                 font=\tiny\sffamily, align=center},
    arr/.style={-{Stealth[length=4pt,width=4pt]}, thick},
    lbl/.style={font=\scriptsize\sffamily, align=center}
  ]

  %--- Hive-on-MapReduce (top) ---
  \node[font=\small\bfseries\sffamily, text=DERed] at (1.4, 7.7) {Hive on MapReduce};
  \node[mr] (m1) at (0,   7.0) {Map 1};
  \node[mr] (r1) at (2.2, 7.0) {Reduce 1};
  \node[disk] (d1) at (3.7, 7.0) {HDFS};
  \node[mr] (m2) at (5.2, 7.0) {Map 2};
  \node[mr] (r2) at (7.4, 7.0) {Reduce 2};
  \node[disk] (d2) at (8.9, 7.0) {HDFS};
  \node[mr] (m3) at (10.4, 7.0) {Map 3};
  \node[mr] (r3) at (12.0, 7.0) {Reduce 3};

  \draw[arr, DERed!60] (m1.east) -- (r1.west);
  \draw[arr, DERed!60] (r1.east) -- (d1.west);
  \draw[arr, DERed!60] (d1.east) -- (m2.west);
  \draw[arr, DERed!60] (m2.east) -- (r2.west);
  \draw[arr, DERed!60] (r2.east) -- (d2.west);
  \draw[arr, DERed!60] (d2.east) -- (m3.west);
  \draw[arr, DERed!60] (m3.east) -- (r3.west);

  \node[lbl, text=DERed] at (3.7, 6.35) {\tiny full write/read};
  \node[lbl, text=DERed] at (8.9, 6.35) {\tiny full write/read};

  %--- Hive on Tez (bottom) ---
  \node[font=\small\bfseries\sffamily, text=DEGreen!70!black] at (4.4, 5.6) {Hive on Tez};

  \node[tez] (v1) at (0,   4.9) {Vertex 1\\(Scan)};
  \node[tez] (v2) at (2.4, 5.3) {Vertex 2\\(Join)};
  \node[tez] (v3) at (2.4, 4.5) {Vertex 3\\(Scan)};
  \node[tez] (v4) at (5.0, 4.9) {Vertex 4\\(Filter)};
  \node[tez] (v5) at (7.6, 4.9) {Vertex 5\\(Aggregate)};
  \node[tez] (v6) at (10.2, 4.9) {Vertex 6\\(Sort)};

  \draw[arr, DEGreen!70!black] (v1.east) -- (v2.west);
  \draw[arr, DEGreen!70!black] (v3.east) -- (v2.west);
  \draw[arr, DEGreen!70!black] (v2.east) -- (v4.west);
  \draw[arr, DEGreen!70!black] (v4.east) -- (v5.west);
  \draw[arr, DEGreen!70!black] (v5.east) -- (v6.west);

  \node[lbl, text=DEGreen!70!black] at (2.4, 3.8)
       {\tiny pipelined: no HDFS writes};
  \draw[DEGreen!50, dashed] (1.5, 4.0) -- (1.9, 4.6);

  % Legend
  \draw[DERed!60, thick] (0.4, 3.1) -- (1.2, 3.1);
  \node[lbl, text=DERed] at (2.6, 3.1) {Hive/MapReduce (forced HDFS materialisation)};
  \draw[DEGreen!70!black, thick] (0.4, 2.7) -- (1.2, 2.7);
  \node[lbl, text=DEGreen!70!black] at (2.3, 2.7) {Hive/Tez (pipelined in memory)};

\end{tikzpicture}
\caption{Hive-on-MapReduce vs.\ Hive-on-Tez for a three-stage query.
         MapReduce forces a full HDFS write and read at each stage boundary
         (shown as cylinders). Tez represents the same query as a DAG of vertices
         connected by pipelined edges; intermediate results flow in memory
         without HDFS materialisation.}
\label{fig:tez-vs-mr}
\end{figure}

Tez introduces three types of edges between vertices that control how data
flows between them:

\begin{itemize}
  \item \textbf{One-to-one edge:} each task of the downstream vertex reads
        from exactly one task of the upstream vertex---no data movement
        required. Used for pipelined filtering and projection.
  \item \textbf{Broadcast edge:} all tasks of the downstream vertex receive
        a full copy of all data produced by the upstream vertex. Used for
        broadcasting a small dimension table to all join tasks (equivalent
        to Spark's broadcast hash join).
  \item \textbf{Scatter-gather edge:} the standard shuffle---each upstream
        task partitions its output by key, and each downstream task reads
        one partition from every upstream task. Used for \texttt{GROUP BY}
        and sort-merge joins.
\end{itemize}

The practical speedup of Hive-on-Tez over Hive-on-MapReduce is 2--5$\times$
for simple queries and up to 100$\times$ for complex multi-stage pipelines,
because the HDFS write-read cost of each stage boundary is eliminated.
By 2015, most Hadoop distributions had configured Hive-on-Tez as the
default execution engine.

%--------------------------------------------------------------------
\subsection{Hive LLAP: In-Memory Interactive Queries}
\label{subsec:llap}
%--------------------------------------------------------------------

Even with Tez, a Hive query incurs several sources of latency that prevent
sub-minute interactive response times. First, each YARN container must be
launched fresh for every query---a process that takes 5--30 seconds of
JVM startup and class-loading overhead. Second, every query must read its
input data from HDFS, even if a prior query just read the same data seconds
ago. Third, thread-level parallelism within a YARN container is limited
to the cores allocated to that container, preventing fine-grained sharing
of executor resources between concurrent queries.

\term{Hive LLAP} (Live Long And Process), introduced in Hive 2.0 (2016),
resolves all three bottlenecks through a fundamentally different execution
model.

\begin{defbox}{Hive LLAP Architecture}
LLAP replaces per-query YARN containers with \textbf{persistent long-running
daemon processes} running on each worker node. Each daemon:
\begin{enumerate}
  \item Maintains a \textbf{columnar in-memory cache} of recently read
        ORC/Parquet data, organised at the column chunk granularity.
  \item Runs a \textbf{shared thread pool} that executes tasks from
        multiple concurrent queries simultaneously, without per-query
        JVM startup overhead.
  \item Performs \textbf{vectorized execution}: processes data in batches
        of 1{,}024 rows using tight CPU loops rather than the
        row-at-a-time interpreted execution model.
\end{enumerate}
\end{defbox}

The in-memory cache is the key feature that enables interactive latency.
A workload in which analysts repeatedly query different slices of the
same large table---typical in a business intelligence environment---will
find that after the first query warms the cache, subsequent queries
across the same table read from RAM at 40--100\GB/s rather than from
HDFS at 1--5\GB/s. Combined with eliminated container startup overhead,
LLAP delivers sub-second latency for queries over cached data, and
5--30 second latency for queries requiring HDFS reads.

%%====================================================================
\section{Query Optimisation in Hive}
\label{sec:hive-optimisation}
%%====================================================================

%--------------------------------------------------------------------
\subsection{Rule-Based vs.\ Cost-Based Optimisation}
\label{subsec:rbo-cbo}
%--------------------------------------------------------------------

In Chapter~\ref{ch:mapreduce}, the Hive compilation pipeline applied
\emph{rule-based optimisation} (RBO): fixed logical rules such as
``push predicates below joins'' and ``eliminate unused columns.'' These
rules improve query plans in many cases but are blind to data
characteristics---they cannot know whether a join's left input has
100 rows or 100 billion rows.

\term{Cost-Based Optimisation} (CBO), introduced in Hive 0.14 (2014)
using Apache Calcite as the optimisation framework, extends the
compilation pipeline with statistics-driven decisions. The CBO reads
table and column statistics from the Hive Metastore (collected by
\texttt{ANALYZE TABLE} commands or auto-gathered on partition creation)
and uses them to estimate the cardinality and selectivity of each
operator in the logical plan. With these estimates, the CBO can:

\begin{itemize}
  \item \textbf{Join reordering}: place the smallest table as the build
        side of a hash join, minimising the memory footprint and build
        time. For a three-table join, the optimal ordering reduces hash
        table size from the cross product of all tables to the smallest
        table's size.
  \item \textbf{Join strategy selection}: choose between Common Join
        (sort-merge, for large tables), Map Join (broadcast, for small
        tables), and Sort-Merge Bucket Join (for pre-bucketed tables)
        based on estimated sizes.
  \item \textbf{Partition pruning with statistics}: estimate the fraction
        of rows eliminated by a predicate and use this to decide whether
        to perform a partition scan or a full-table scan.
\end{itemize}

Enabling the CBO requires first collecting statistics:

\begin{lstlisting}[style=hiveql, caption={Collecting Hive statistics for CBO}]
-- Collect table-level and column-level statistics
ANALYZE TABLE orders COMPUTE STATISTICS;
ANALYZE TABLE orders COMPUTE STATISTICS FOR COLUMNS
    customer_id, order_date, amount, region;

-- Verify statistics in the Metastore
DESCRIBE FORMATTED orders;
-- Output includes: numRows, numFiles, totalSize,
--                  column-level: min, max, numNulls, numDVs, avgLen
\end{lstlisting}

%--------------------------------------------------------------------
\subsection{Join Execution Strategies}
\label{subsec:join-strategies}
%--------------------------------------------------------------------

Hive implements four distinct join strategies, each suited to different
combinations of table sizes and data organisation. Selecting the wrong
join strategy is one of the most common causes of avoidable performance
degradation in Hive ETL pipelines.

\textbf{Common Join} (Reduce-Side Join): the default for large tables.
Both tables are shuffled on the join key to co-locate matching rows on
the same reducer. For two tables of size $A$ and $B$, the shuffle
transfers $A + B$ bytes of data across the network. The shuffle is the
bottleneck: for terabyte-scale tables, it can require hours.

\textbf{Map Join} (Broadcast Join): applicable when one table is small
enough to fit in executor memory (default threshold: 25\MB\ in Hive,
10\MB\ in Spark SQL). The small table is loaded into a hash table in
memory at each Map task. The large table is scanned without any shuffle:
each Map task probes its local hash table to find matching rows. No
reducer is needed. For a join of a 10\TB\ fact table against a 20\MB\
dimension table, Map Join eliminates the entire 10\TB\ shuffle.

\textbf{Bucket Map Join}: applicable when both tables are bucketed on
the join key with the same number of buckets (or one is a multiple of
the other). Each Map task processes only the corresponding bucket of each
table, without loading the entire small table. Extends Map Join to
cases where the ``small'' table is too large for a full broadcast but
is well-organised.

\textbf{Sort-Merge Bucket Join} (SMB Join): the highest-performance join
for large tables that are both bucketed \emph{and} sorted on the join key.
Because matching rows in both tables are co-located (same bucket) and
pre-sorted, the join can be performed by a single linear scan of each
bucket pair---analogous to the merge step in merge sort---without any
shuffle whatsoever. SMB Join requires upfront investment in table design
(bucketing and sorting at write time) but delivers the best possible
join performance at read time.

\begin{table}[ht]
\centering
\caption{Hive join strategy comparison}
\label{tab:join-strategies}
\renewcommand{\arraystretch}{1.6}
\begin{tabular}{L{2.4cm} L{1.6cm} L{1.6cm} L{4.6cm}}
\toprule
\textbf{Strategy} & \textbf{Shuffle?} & \textbf{Requirement} & \textbf{Optimal use case} \\
\midrule
Common Join      & Yes (full) & None          & Large–large join; no pre-processing \\
Map Join         & No         & Small table $<$ threshold & Large table $\times$ small dimension \\
Bucket Map Join  & No         & Both bucketed on join key & Medium–large join; data pre-bucketed \\
SMB Join         & No         & Both bucketed + sorted    & Largest tables; highest write cost \\
\bottomrule
\end{tabular}
\end{table}

%--------------------------------------------------------------------
\subsection{Vectorized Execution}
\label{subsec:vectorized}
%--------------------------------------------------------------------

Traditional SQL execution engines process one row at a time through
a tree of operator nodes---the \term{Volcano model} (or pull-based
model). For each row, the query engine invokes virtual function calls
at each operator, navigating the operator tree once per record. At
billions of rows, the overhead of these function calls and the
corresponding CPU branch mispredictions becomes significant.

\term{Vectorized execution}, enabled in Hive with
\texttt{SET hive.vectorized.execution.enabled = true}, processes data
in batches of 1{,}024 rows (one \emph{vector batch}) rather than one
row at a time. Each operator receives a batch of column vectors and
applies a tight loop over all 1{,}024 rows without virtual function
calls. The CPU can predict the loop branch perfectly and pipeline the
arithmetic within each column chunk, making effective use of SIMD
(Single Instruction Multiple Data) instructions on modern hardware.

For CPU-bound operations (arithmetic on numerical columns, string
comparisons, hash computations), vectorized execution provides a
2--10$\times$ throughput improvement over the row-at-a-time model.
ORC files (Chapter~\ref{ch:yarn-formats}) are particularly well-suited
for vectorized execution because their columnar layout naturally delivers
data to the engine in column-major order, matching the layout expected
by the vector batch operators.

%%====================================================================
\section{Presto and Trino: SQL on Everything}
\label{sec:presto}
%%====================================================================

\term{Presto} was created at Facebook in 2012 to solve a specific problem:
the Facebook data warehouse team had hundreds of analysts running queries
on petabyte-scale data, and Hive-on-MapReduce was too slow for the
interactive ad-hoc analytical queries that characterise analyst workflows.
A query that answers the question ``how many users in Germany clicked on
ad campaign 4421 in the last 7 days?'' should return in seconds, not hours.

After open-sourcing Presto in 2013, it was adopted by hundreds of
organisations including LinkedIn, Twitter, Airbnb, and Uber.
In 2019, the original Presto creators left Facebook and forked the project
as \term{Trino} (2020). Both Presto (now maintained by the Presto
Foundation) and Trino remain actively developed, with largely compatible
SQL dialects. In this chapter, ``Presto'' refers to both unless the
distinction matters.

\parencite{Presto2019}

%--------------------------------------------------------------------
\subsection{Presto Architecture}
\label{subsec:presto-arch}
%--------------------------------------------------------------------

A Presto cluster consists of one \term{Coordinator} node and multiple
\term{Worker} nodes, communicating over HTTP/2.

\begin{figure}[ht]
\centering
\begin{tikzpicture}[
    coord/.style={rectangle, draw=DEBlue!70, fill=DELightBlue!50,
                  rounded corners=5pt, minimum width=4.8cm, minimum height=0.9cm,
                  font=\small\bfseries\sffamily, align=center},
    subcomp/.style={rectangle, draw=DEBlue!40, fill=DEBlue!8,
                    minimum width=2.0cm, minimum height=0.6cm,
                    font=\tiny\sffamily, align=center},
    worker/.style={rectangle, draw=DEGreen!70, fill=DELightGreen,
                   rounded corners=4pt, minimum width=2.4cm, minimum height=0.7cm,
                   font=\small\sffamily, align=center},
    conn/.style={rectangle, draw=DEAccent!70, fill=DELightAccent,
                 rounded corners=3pt, minimum width=1.8cm, minimum height=0.55cm,
                 font=\tiny\sffamily, align=center},
    client/.style={rectangle, draw=DEGray!60, fill=DEGray!15,
                   rounded corners=3pt, minimum width=1.8cm, minimum height=0.55cm,
                   font=\small\sffamily, align=center},
    arr/.style={-{Stealth[length=4pt,width=4pt]}, thick},
    lbl/.style={font=\scriptsize\sffamily, align=center}
  ]

  % Client
  \node[client] (cli) at (0, 8.8) {Client\\(JDBC/REST)};

  % Coordinator box
  \node[coord] (coord) at (4.5, 8.8) {Coordinator};
  \node[subcomp] (parser)  at (3.2, 7.8) {Parser};
  \node[subcomp] (planner) at (5.8, 7.8) {Query Planner};
  \node[subcomp] (sched)   at (3.2, 7.1) {Scheduler};
  \node[subcomp] (meta)    at (5.8, 7.1) {Metadata (Catalogs)};

  \draw[arr, DEBlue!60] (cli.east) -- (coord.west)
       node[midway, above, lbl] {\tiny SQL query};
  \draw[arr, DEBlue!30] (coord.south) -- (parser.north);
  \draw[arr, DEBlue!30] (coord.south) -- (planner.north);
  \draw[arr, DEBlue!30] (parser.south) -- (sched.north);
  \draw[arr, DEBlue!30] (planner.south) -- (meta.north);

  % Workers
  \node[worker] (w1) at (0.5,  5.6) {Worker 1};
  \node[worker] (w2) at (3.0,  5.6) {Worker 2};
  \node[worker] (w3) at (5.5,  5.6) {Worker 3};
  \node[worker] (w4) at (8.0,  5.6) {Worker 4};

  \draw[arr, DEBlue] (sched.south) to[out=-90, in=90] (w1.north);
  \draw[arr, DEBlue] (sched.south) to[out=-90, in=90] (w2.north);
  \draw[arr, DEBlue] (sched.south) to[out=-90, in=90] (w3.north);
  \draw[arr, DEBlue] (sched.south) to[out=-90, in=90] (w4.north);
  \node[lbl, text=DEBlue] at (4.5, 6.5) {\tiny distribute tasks};

  % Workers pipeline between themselves
  \draw[arr, DEGreen!70!black, dashed]
       (w1.east) -- (w2.west)
       node[midway, above, lbl, text=DEGreen!70!black] {\tiny pipeline};
  \draw[arr, DEGreen!70!black, dashed] (w2.east) -- (w3.west);
  \draw[arr, DEGreen!70!black, dashed] (w3.east) -- (w4.west);

  % Data sources (connectors)
  \node[conn] (c1) at (0.5,  4.3) {Hive\\(HDFS/S3)};
  \node[conn] (c2) at (3.0,  4.3) {JDBC\\(MySQL)};
  \node[conn] (c3) at (5.5,  4.3) {Kafka};
  \node[conn] (c4) at (8.0,  4.3) {Elastic-\\search};

  \draw[arr, DEAccent] (w1.south) -- (c1.north);
  \draw[arr, DEAccent] (w2.south) -- (c2.north);
  \draw[arr, DEAccent] (w3.south) -- (c3.north);
  \draw[arr, DEAccent] (w4.south) -- (c4.north);
  \node[lbl, text=DEAccent] at (4.5, 3.7) {\tiny Connectors (read data)};

\end{tikzpicture}
\caption{Presto cluster architecture. The Coordinator receives SQL queries from clients,
         parses and plans them, then distributes tasks to Worker nodes. Workers
         pipeline results between themselves in memory. Each worker reads data
         through a Connector that abstracts the underlying data source---HDFS,
         JDBC databases, Kafka, Elasticsearch, or others.}
\label{fig:presto-arch}
\end{figure}

The \term{Coordinator} is the query planning and scheduling process:
\begin{itemize}
  \item \textbf{Parser}: converts the SQL string into an abstract syntax
        tree.
  \item \textbf{Analyser}: resolves table names and column types by
        querying the registered catalogs (each catalog is backed by a
        connector's metadata interface).
  \item \textbf{Query Planner}: produces a distributed execution plan
        as a tree of \emph{plan fragments}. Each fragment corresponds to
        a set of tasks that can run concurrently on multiple workers.
        Plan fragments are connected by \emph{exchange operators}: the
        upstream fragment's output is buffered in the upstream workers'
        memory and pulled by the downstream fragment over HTTP.
  \item \textbf{Scheduler}: assigns fragments to specific workers, taking
        data locality (HDFS block locations) into account for the scan
        fragments.
\end{itemize}

\term{Workers} execute the tasks assigned by the coordinator. Each worker
processes a subset of the data by reading through its connector and
applying the plan fragment's operators (filter, project, aggregate, join).
Workers communicate intermediate results to each other via in-memory
buffers over HTTP/2. \emph{There is no disk I/O between plan fragments}:
all intermediate data flows in memory. This is the primary architectural
difference from Hive-on-MapReduce.

%--------------------------------------------------------------------
\subsection{The Connector Model: SQL on Everything}
\label{subsec:connectors}
%--------------------------------------------------------------------

The most distinctive feature of Presto is its \term{connector model}:
a pluggable interface that allows Presto to read from (and in some cases
write to) any data source that implements the connector API. Each connector
provides:
\begin{itemize}
  \item \textbf{Metadata API:} listing catalogs, schemas, tables, columns,
        and column statistics.
  \item \textbf{Data location API:} describing how data is split into
        independently-readable chunks (HDFS splits, Kafka partition
        offsets, JDBC pagination ranges).
  \item \textbf{Record cursor / Page source API:} reading actual data
        records from the underlying storage.
\end{itemize}

This abstraction enables \term{federated queries}---SQL queries that
join data from multiple different data sources in a single statement:

\begin{lstlisting}[style=hiveql, caption={Presto federated query: joining HDFS data with a MySQL table}]
-- Join HDFS order data (Hive connector)
-- with a MySQL reference table (JDBC connector)
-- in a single Presto query
SELECT  h.customer_id,
        m.customer_name,
        m.country,
        SUM(h.order_amount)  AS total_spent,
        COUNT(*)             AS order_count
FROM    hive.analytics.orders  AS h          -- HDFS/Parquet via Hive connector
JOIN    mysql.crm.customers    AS m          -- MySQL via JDBC connector
     ON h.customer_id = m.customer_id
WHERE   h.order_date >= DATE '2024-01-01'
  AND   m.country IN ('US', 'GB', 'DE')
GROUP BY h.customer_id, m.customer_name, m.country
ORDER BY total_spent DESC
LIMIT 100;
\end{lstlisting}

This query reads the \texttt{orders} table from Hive (backed by Parquet
files on HDFS) and the \texttt{customers} table from a MySQL database,
joining them in Presto's memory without any ETL pipeline or intermediate
data movement. In a Hive-only environment, this query would require first
exporting the MySQL table to HDFS and creating a Hive external table---a
process taking hours.

%--------------------------------------------------------------------
\subsection{Presto vs.\ Hive: Design Philosophy}
\label{subsec:presto-vs-hive}
%--------------------------------------------------------------------

The fundamental design difference between Presto and Hive reflects a
deliberate trade-off between \emph{query latency} and \emph{fault tolerance}.

\textbf{Presto's design choices:}
\begin{itemize}
  \item \textbf{No intermediate disk spill (by default).} All data
        processing is in-memory. A query that exceeds the cluster's
        total memory fails immediately with an ``exceeded memory limit''
        error rather than spilling to disk and completing slowly.
  \item \textbf{No task retry.} If a worker node fails during query
        execution, the entire query fails. The coordinator does not
        retry failed tasks on different workers (though more recent
        Presto/Trino versions have added limited retry capabilities).
  \item \textbf{Low startup overhead.} Workers are persistent processes;
        there is no container-launch overhead. A query begins executing
        within milliseconds of submission.
  \item \textbf{Push-based pipelining.} Results flow from upstream to
        downstream operators without waiting for the upstream stage to
        complete---the first results are produced before the last input
        row is read.
\end{itemize}

\textbf{Hive's design choices:}
\begin{itemize}
  \item \textbf{Disk-based fault tolerance.} Intermediate results are
        materialised to HDFS at each stage boundary (MapReduce) or
        can be spilled to local disk (Tez). A task can be re-executed
        if it fails, without restarting the entire query.
  \item \textbf{Unbounded query size.} A Hive query can process any
        volume of data, regardless of cluster RAM, because it spills
        to HDFS. A 100\TB\ sort that would exhaust Presto's memory
        completes in Hive---slowly, but completely.
  \item \textbf{Higher startup overhead.} YARN container launches add
        5--30 seconds to every query. LLAP partially mitigates this
        with persistent daemons.
\end{itemize}

\begin{caution}{Presto is Not a Hive Replacement for ETL}
A common mistake is deploying Presto as a general-purpose replacement
for Hive and then running large ETL pipelines through it. A 50\TB\
nightly ETL pipeline that produces many intermediate materialised results
belongs in Hive-on-Tez (or Spark), where disk-based fault tolerance
and HDFS spill prevent out-of-memory failures. Presto is optimised for
the interactive analytics use case: queries that return in seconds or
minutes, over data that fits in the cluster's RAM. Running multi-hour
ETL jobs through Presto risks losing all progress if any worker fails.
Use the right engine for each workload.
\end{caution}

%%====================================================================
\section{The SQL-on-Hadoop Ecosystem}
\label{sec:sql-ecosystem}
%%====================================================================

Hive and Presto are the two most widely deployed SQL-on-Hadoop engines,
but the ecosystem includes several other significant systems. Understanding
their design points helps data engineers match engine to workload.

%--------------------------------------------------------------------
\subsection{Apache Impala}
\label{subsec:impala}
%--------------------------------------------------------------------

\term{Apache Impala}, developed by Cloudera and released in 2012, is a
massively parallel processing (MPP) SQL engine designed for interactive
analytics. Unlike Hive, which is built on top of MapReduce/Tez, Impala
is implemented entirely in C++ and runs as a native process on each
cluster node---bypassing the JVM and YARN container model entirely.

Impala's key design characteristics:
\begin{itemize}
  \item \textbf{No JVM.} C++ avoids JVM garbage collection pauses,
        which can add hundreds of milliseconds of latency to individual
        query tasks---unacceptable for sub-second interactive queries.
  \item \textbf{Shared metadata.} Impala reads the Hive Metastore for
        table definitions, enabling analysts to switch between Hive and
        Impala queries against the same tables.
  \item \textbf{In-memory, no fault tolerance.} Like Presto, Impala
        does not retry failed tasks. Queries that exceed memory fail.
  \item \textbf{HDFS and Kudu.} Impala supports both HDFS (Parquet, ORC)
        and Apache Kudu (a columnar storage engine optimised for fast
        random reads and updates), making it suitable for near-real-time
        analytics over mutable data.
\end{itemize}

%--------------------------------------------------------------------
\subsection{Comparative Analysis}
\label{subsec:engine-comparison}
%%--------------------------------------------------------------------

\begin{table}[ht]
\centering
\caption{SQL-on-Hadoop engine comparison}
\label{tab:engine-comparison}
\renewcommand{\arraystretch}{1.6}
\begin{tabular}{L{1.8cm} L{1.7cm} L{1.5cm} L{1.8cm} L{2.0cm} L{2.2cm}}
\toprule
\textbf{Engine} & \textbf{Latency} & \textbf{Scale} & \textbf{Fault Tol.} & \textbf{Data sources} & \textbf{Best use case} \\
\midrule
Hive/MapReduce  & Hours     & Unlimited & Full       & HDFS (ORC/Parquet/CSV) & Multi-TB batch ETL \\
Hive/Tez        & Minutes   & Unlimited & Full       & HDFS                   & Complex batch analytics \\
Hive/LLAP       & Seconds   & Large     & Full       & ORC/Parquet on HDFS    & BI dashboards; repeated queries \\
Presto/Trino    & Seconds   & RAM-bound & Query-fail & Federated (Hive, JDBC, Kafka, \ldots) & Ad-hoc interactive analytics \\
Spark SQL       & Seconds   & Large     & Partial    & HDFS, S3, JDBC, Kafka  & ML pipelines; unified batch+stream \\
Apache Impala   & Sub-second & RAM-bound & Query-fail & HDFS, Kudu             & Fastest interactive on Parquet \\
\bottomrule
\end{tabular}
\end{table}

The selection principle is workload-driven:
\begin{itemize}
  \item \textbf{Nightly batch ETL over terabytes:} Hive-on-Tez or Spark.
        Fault tolerance and unlimited scale matter; latency does not.
  \item \textbf{BI dashboards with repeated queries:} Hive LLAP.
        The cache warms on the first query; subsequent queries are served
        from memory at sub-second latency.
  \item \textbf{Ad-hoc analyst queries ($<$5 minutes):} Presto/Trino.
        Federated query capability and low startup latency are critical.
  \item \textbf{Fastest interactive queries on Parquet, low cardinality:}
        Apache Impala. C++ execution with no JVM overhead.
  \item \textbf{ML training on the same data as SQL analytics:} Spark SQL.
        A single framework covers both SQL and MLlib.
\end{itemize}

\begin{researchspotlight}{Venkataraman et al.\ (2019) --- Presto: SQL on Everything}
\textbf{Citation:} Sethi, R., Traverso, M., Sundstrom, D., Phillips, D.,
Xie, W., Sun, Y., Yigitbasi, N., Jin, H., Hwang, E., Shingte, N., and
Berner, C.\ (2019). Presto: SQL on Everything. \textit{Proceedings of the
35th IEEE International Conference on Data Engineering (ICDE 2019)},
pp.\ 1802--1813. \parencite{Presto2019}

\smallskip\noindent
\textbf{Key insight:} The paper's thesis is that the connector model---a
clean separation between the SQL execution engine and the data source---
is the enabling abstraction for ``SQL on everything.'' By standardising
how the engine asks a data source about its schema and data layout, Presto
can treat any data source (HDFS, RDBMS, Kafka, Elasticsearch, custom
REST APIs) identically from the query execution perspective. The paper
calls this property \emph{data source agnosticism}.

\smallskip\noindent
\textbf{Technical contributions:}
\begin{itemize}
  \item Described the Coordinator/Worker/Connector architecture and
        the push-based distributed execution model, showing how
        pipelined query plans eliminate the HDFS materialisation
        bottleneck of Hive-on-MapReduce.
  \item Introduced the \emph{Connector SPI} (Service Provider Interface)---
        the five interfaces (metadata, data location, data source, record
        cursor, and transaction) that a connector must implement to integrate
        a new data source into Presto without modifying the core engine.
  \item Reported Facebook's production deployment: Presto ran on thousands
        of worker nodes, processing petabytes per day, serving 30{,}000+
        queries per day from 1{,}000+ internal Facebook users with a
        median query latency of under 1 minute and a P95 latency of under
        10 minutes.
\end{itemize}

\smallskip\noindent
\textbf{Impact today:} Presto/Trino is deployed at hundreds of organisations
and has become the de facto standard engine for interactive SQL analytics
on cloud data lakes. Amazon Athena (serverless SQL on S3) is built on
Presto. Databricks SQL Warehouse and Snowflake adopted similar connector
architectures. The Presto SPI was the direct inspiration for Apache
Arrow's DataFusion connector model and the Delta Sharing protocol's
reader interface.
\end{researchspotlight}

%%====================================================================
\section{Case Study: Airbnb's Migration to a Unified SQL Platform}
\label{sec:airbnb-case}
%%====================================================================

\begin{casestudy}{Airbnb --- From Hive Monolith to Hive-Plus-Presto Tiered SQL Platform}

\textbf{Background.}
By 2015, Airbnb had grown to over 2 million listings and 40 million
users, with a data team of approximately 100 engineers and analysts
generating more than 10{,}000 Hive queries per day. The company's
data warehouse ran on a 500-node Hadoop cluster with approximately
15\PB\ of data, processed exclusively by Hive-on-MapReduce.

The data platform was suffering from three interconnected problems:
analyst query latency averaged 45 minutes (with some queries taking
8+ hours), cluster utilisation was chronically at 95\%+ because
long-running batch ETL jobs and short analyst queries shared the same
YARN queue, and analyst adoption of data-driven decision-making was
limited because the 45-minute query latency discouraged exploratory
iteration.

\smallskip
\textbf{Phase 1 (2016): Hive Tez Migration.}
Airbnb's first performance initiative migrated all analyst-facing Hive
queries from MapReduce to Tez. The migration was transparent to analysts:
no query rewrites were required, as the HiveQL syntax is identical
regardless of execution engine. The results were significant:

\begin{itemize}
  \item Median analyst query latency dropped from 45 minutes to 18 minutes.
  \item P95 latency dropped from $>$8 hours to 75 minutes.
  \item Cluster CPU utilisation dropped from 95\% to 70\% for the same
        analytical workload, because Tez eliminated the HDFS
        materialisation I/O that had been consuming a disproportionate
        share of disk bandwidth.
\end{itemize}

However, even 18 minutes remained too slow for the exploratory, iterative
analysis style that Airbnb's product and growth teams required. An analyst
testing hypotheses about booking conversion rates might run 15--20
queries in sequence; at 18 minutes each, this required a full day.

\smallskip
\textbf{Phase 2 (2017): Presto Deployment for Interactive Analytics.}
Airbnb deployed a dedicated 150-node Presto cluster alongside the
existing Hive cluster. The Presto cluster shared the Hive Metastore,
so all tables defined for Hive were immediately queryable through Presto
without any schema migration.

The workload routing strategy divided queries by expected duration:
\begin{center}
\begin{tabular}{lll}
\toprule
\textbf{Expected duration} & \textbf{Engine} & \textbf{Queue} \\
\midrule
$<$5 minutes (interactive) & Presto & \texttt{analyst\_interactive} \\
5--60 minutes (analytical) & Hive/Tez & \texttt{analyst\_batch} \\
$>$60 minutes (ETL/batch)  & Hive/Tez & \texttt{production\_etl} \\
\bottomrule
\end{tabular}
\end{center}

Routing was implemented in Airbnb's internal query routing service,
which classified incoming queries based on a set of heuristics---estimated
data scanned (from partition statistics), presence of \texttt{CROSS JOIN}
or unconstrained \texttt{JOIN} patterns, and explicit analyst preference
flags. The routing service directed queries to the appropriate cluster
transparently.

\smallskip
\textbf{The booking funnel analysis pipeline.}
A representative workload that benefited most from Presto was the
booking funnel analysis: understanding how users progress from search
to booking across Airbnb's platform. The analysis requires joining four
tables:

\begin{enumerate}
  \item \textbf{\texttt{search\_events}} ($\sim$2\TB/day): one row
        per guest search, with search parameters, result count, and
        filters applied.
  \item \textbf{\texttt{listing\_views}} ($\sim$500\GB/day): one row
        per listing viewed from a search results page.
  \item \textbf{\texttt{booking\_attempts}} ($\sim$50\GB/day): one row
        per booking initiation, including payment method and price.
  \item \textbf{\texttt{dim\_listings}} ($\sim$4\GB, daily snapshot):
        listing metadata including location, price tier, and amenity flags.
\end{enumerate}

The booking funnel query computes conversion rates at each step of the
funnel, segmented by geography and listing price tier. On Hive-on-MapReduce,
this query ran in 4 hours. On Hive-on-Tez, it ran in 55 minutes. On
Presto, after Catalyst-equivalent optimisations by the Presto analyser
(broadcast join for \texttt{dim\_listings}, which at 4\GB\ fits within
each worker's broadcast buffer), the query ran in 4 minutes.

The 4-minute latency transformed analyst behaviour: instead of running
the funnel analysis once per day and committing to the results, analysts
began running it hourly during A/B test periods, checking whether a new
feature was affecting conversion rates in near real time.

\smallskip
\textbf{Phase 3 (2019): LLAP for Dashboards.}
Airbnb's engineering dashboards required even lower latency for
recurring queries: the homepage executive dashboard refreshed every
5 minutes, pulling the same suite of KPIs (bookings, revenue, host
activations, guest satisfaction scores) from Hive tables. On Presto,
these queries took 20--60 seconds depending on cluster load.

Deploying Hive LLAP for the dashboard tier cached the relevant table
partitions in the LLAP daemons' memory after the first query of each
5-minute cycle. Subsequent queries in the same cycle served from the
in-memory cache, reducing dashboard query latency to 1--5 seconds.

\smallskip
\textbf{Outcomes.}
The three-tier SQL architecture---LLAP for dashboards, Presto for
interactive analytics, Hive/Tez for batch ETL---produced measurable
improvements across every operational dimension:

\begin{itemize}
  \item \textbf{Analyst query latency:} P50 reduced from 45 minutes
        (Hive/MR) to 3.5 minutes (Presto), a $13\times$ improvement.
  \item \textbf{Data-driven decisions:} analyst-submitted queries per day
        increased from 10{,}000 to 47{,}000 within 18 months of
        deploying Presto---a $4.7\times$ increase, attributed to the
        lower iteration cost of interactive query response times.
  \item \textbf{Cluster efficiency:} separating interactive (Presto)
        from batch ETL (Hive/Tez) eliminated the queue starvation
        problem. Long-running ETL jobs no longer delayed analyst queries;
        the Presto cluster's dedicated allocation guaranteed
        sub-5-minute response for all routed queries.
  \item \textbf{Engineering savings:} the number of ``bespoke
        pre-computed aggregate tables'' (summary tables created by
        engineers specifically to speed up slow analyst queries)
        dropped by 60\%, because analysts could run the original
        queries directly on Presto at acceptable latency.
\end{itemize}

\textbf{The architectural lesson.}
Airbnb's platform evolution illustrates the principle that no single
SQL engine is optimal for all workloads. The cost of this differentiation---
operating three distinct query engines against the same Hive Metastore---
is engineering complexity in routing, monitoring, and SLA management.
The benefit is that each engine operates within its optimal design envelope:
LLAP for cached sub-second BI, Presto for interactive exploration, Hive/Tez
for batch ETL that requires fault tolerance and unlimited scale.
\end{casestudy}

%%====================================================================
\section*{Chapter Summary}
\label{sec:ch7-summary}
%%====================================================================

\begin{itemize}
  \item HiveQL's \textbf{complex data types} (\texttt{ARRAY}, \texttt{MAP},
        \texttt{STRUCT}) model semi-structured data natively.
        \textbf{LATERAL VIEW EXPLODE} unnests array and map columns into
        relational rows, enabling standard \texttt{GROUP BY} aggregation
        over nested data without UDFs.

  \item \textbf{Window functions} (\texttt{OVER}, \texttt{PARTITION BY},
        \texttt{ORDER BY}, \texttt{ROWS BETWEEN}) compute running totals,
        moving averages, and rankings without collapsing the result set.
        Key functions: \texttt{ROW\_NUMBER}, \texttt{RANK},
        \texttt{LAG}, \texttt{LEAD}, \texttt{NTILE}.

  \item Hive \textbf{UDF} categories: scalar UDF (one-to-one),
        \textbf{UDAF} (many-to-one, must implement partial + final phases),
        and \textbf{UDTF} (one-to-many, used for \texttt{EXPLODE}-style
        expansion).

  \item \textbf{Apache Tez} replaces MapReduce's chain of jobs with a
        single DAG of vertices, eliminating forced HDFS materialisation
        at each stage boundary. Speedup: 2--5$\times$ for simple queries,
        up to 100$\times$ for complex multi-stage pipelines.

  \item \textbf{Hive LLAP} adds a persistent in-memory columnar cache
        and shared executor threads, eliminating per-query JVM startup
        overhead and enabling sub-second latency for cached data.

  \item Hive's \textbf{Cost-Based Optimizer} (CBO, powered by Apache
        Calcite) uses Metastore statistics (\texttt{ANALYZE TABLE}) to
        choose join order and join strategy. Without statistics, the CBO
        falls back to rule-based optimisation.

  \item The four Hive \textbf{join strategies}: Common Join (shuffle,
        always applicable), Map Join (no shuffle, small table broadcast),
        Bucket Map Join (no shuffle, both bucketed), and Sort-Merge Bucket
        Join (no shuffle, both bucketed and sorted --- highest performance,
        highest write cost).

  \item \textbf{Vectorized execution} processes 1{,}024 rows per batch
        rather than one row at a time, eliminating virtual function call
        overhead and enabling SIMD parallelism. Compatible with ORC's
        columnar layout. Typical speedup: 2--10$\times$ for CPU-bound
        operations.

  \item \textbf{Presto/Trino} uses a Coordinator + Worker architecture
        with fully pipelined, in-memory execution and no fault tolerance.
        The \textbf{Connector SPI} enables federated SQL across HDFS,
        JDBC databases, Kafka, and arbitrary data sources in a single
        query.

  \item Key Presto vs.\ Hive trade-off: Presto sacrifices fault tolerance
        and unbounded scale for low latency and federated access. Use
        Presto for interactive queries ($<$5 minutes); use Hive for batch
        ETL (hours) that requires fault tolerance and disk spill.
\end{itemize}

%%====================================================================
\section*{Exercises}
\label{sec:ch7-exercises}
%%====================================================================

\exercisesection{Short Questions}

\sq What is the \texttt{LATERAL VIEW EXPLODE} construct in HiveQL?
A table \texttt{orders} has a column \texttt{tags ARRAY<STRING>}.
Write a HiveQL query that counts the total number of times each tag
appears across all orders.

\sq What is the difference between \texttt{RANK()} and
\texttt{DENSE\_RANK()}? Give a concrete example with five rows where the
two functions produce different output.

\sq Explain the three categories of Hive UDFs: UDF, UDAF, and UDTF.
For each, give one real-world use case and explain why it falls into
that category.

\sq What problem does Apache Tez solve that Hive-on-MapReduce does not?
For a HiveQL query that compiles to three MapReduce jobs with 500\MB\
of intermediate data at each stage boundary, calculate the HDFS I/O
saved by using Tez instead.

\sq What is Hive LLAP? Describe the three components of LLAP architecture
and explain why LLAP enables sub-second latency for repeated queries
over the same table.

\sq What is the difference between rule-based optimisation (RBO) and
cost-based optimisation (CBO) in Hive? Give one example where RBO would
produce a suboptimal plan that CBO would avoid.

\sq Explain Hive's Sort-Merge Bucket Join. What preconditions must
both tables satisfy, and why does this join require no shuffle at query
time?

\sq What is the Presto Connector SPI? Name three data sources that
Presto can query through connectors, and explain what interface the
connector must implement for each source.

\sq A data engineer argues that Presto should replace Hive entirely
because Presto is always faster. Write a detailed counter-argument
using a specific workload scenario where Hive is the correct choice.

\sq HiveServer2 replaced the original Hive command-line tool. What
two problems with the original tool did HiveServer2 address, and what
interface (protocol) does HiveServer2 expose to clients?

\sq What is vectorized execution in Hive? Explain the difference from
the Volcano (row-at-a-time) execution model, and identify which file
format is best suited for vectorized execution in Hive.

\sq A Presto query fails with ``Query exceeded maximum memory.''
Explain: (a) why this error occurs; (b) why the same query would
not fail in Hive-on-Tez; (c) what you should do to fix the problem
without switching engines.

\sq Explain Hive's Map Join strategy. What is the default table size
threshold for a Map Join, and what happens during query execution
that eliminates the shuffle?

\sq What is a Presto federated query? Write the SQL statement structure
(not actual data) for a query that joins a Hive table with a MySQL table.
What does the Presto Coordinator need to execute this query?

\sq Compare Apache Impala and Presto on three dimensions: implementation
language, JVM usage, and approach to data source support.

\bigskip
\exercisesection{Analytical Problems}

\ap \textbf{Window Function Design.}
A subscription streaming service wants to analyse user churn. The table
\texttt{user\_daily\_activity} has columns: \texttt{user\_id},
\texttt{activity\_date}, \texttt{minutes\_streamed} (0 if not active),
\texttt{country}.

\begin{enumerate}[label=(\alph*)]
  \item Write a HiveQL window function query that computes, for each
        user and date, the 7-day rolling average of \texttt{minutes\_streamed}.
  \item Write a query using \texttt{LAG} to compute the number of days
        since each user last had a non-zero \texttt{minutes\_streamed}
        value.
  \item Write a query using \texttt{NTILE(10)} to divide users into
        10 deciles by their total minutes streamed in the past 30 days,
        and compute the churn rate for each decile.
  \item Explain how the window function queries in (a)--(c) would be
        compiled by Hive-on-Tez compared to Hive-on-MapReduce. How many
        jobs/stages would each require?
\end{enumerate}

\ap \textbf{Join Strategy Selection.}
A retail analytics system has the following tables:
\begin{center}
\begin{tabular}{llll}
\toprule
\textbf{Table} & \textbf{Rows} & \textbf{Compressed size} & \textbf{Bucketed on} \\
\midrule
\texttt{transactions}  & 80 billion & 12\TB & \texttt{customer\_id} (500 buckets) \\
\texttt{customers}     & 200 million & 80\GB & \texttt{customer\_id} (500 buckets) \\
\texttt{dim\_country}  & 250         & 40\KB & None \\
\texttt{dim\_category} & 8{,}000     & 2\MB  & None \\
\bottomrule
\end{tabular}
\end{center}

\begin{enumerate}[label=(\alph*)]
  \item For a join between \texttt{transactions} and \texttt{dim\_country},
        identify the optimal join strategy and explain why. Estimate the
        network traffic for the alternative (Common Join) approach.
  \item For a join between \texttt{transactions} and \texttt{customers},
        both bucketed on \texttt{customer\_id} with 500 buckets, identify
        the optimal join strategy. What additional condition must hold for
        Sort-Merge Bucket Join to be applicable?
  \item For a three-table join: \texttt{transactions} JOIN
        \texttt{customers} JOIN \texttt{dim\_category}, propose the join
        order and strategy for each join. Justify the order using the CBO's
        cost minimisation objective.
  \item Explain what Hive statistics (output of \texttt{ANALYZE TABLE})
        the CBO would use to automatically determine the join strategies
        in (a)--(c) without manual hints.
\end{enumerate}

\ap \textbf{Presto vs.\ Hive Workload Assignment.}
A data analytics team runs the following five daily workloads:
\begin{enumerate}[label=(\roman*)]
  \item A nightly ETL job that reads 200\TB\ of raw clickstream logs,
        filters and aggregates them, and writes the result to a Hive
        table (runtime: 4--6 hours).
  \item Ad-hoc analyst queries over the last 7 days of clickstream data
        ($\sim$50\GB) to investigate a product anomaly (target: $<$3 min).
  \item A weekly cohort analysis joining three tables totalling 800\GB
        with 12 join stages and complex window functions (runtime target:
        $<$30 minutes).
  \item A live executive dashboard refreshing every 5 minutes, running
        the same 5 KPI queries over the last hour of data ($\sim$2\GB).
  \item A machine learning feature engineering pipeline that generates
        1{,}000 features per user from 60 days of history, followed by
        a logistic regression training run (runtime: 2--3 hours).
\end{enumerate}

For each workload, specify: (a) which engine (Hive/MapReduce, Hive/Tez,
Hive/LLAP, Presto, Spark SQL) is most appropriate and why; (b) the
specific property of the workload that drives the engine choice;
(c) the risk of using the wrong engine.

\ap \textbf{LLAP Cache Sizing.}
An analytics team deploys Hive LLAP on a 50-node cluster. Each node
has 128\GB\ of RAM. The LLAP daemon on each node is allocated 80\GB,
of which 60\GB\ is dedicated to the in-memory columnar cache and 20\GB\
to the executor thread pool.

The team's most frequently queried table is \texttt{sales\_events}:
partitioned by day, each day's partition is 400\GB\ compressed (ORC)
but 1.2\TB\ uncompressed. Analysts typically query the last 30 days.

\begin{enumerate}[label=(\alph*)]
  \item Calculate the total LLAP cache capacity across the cluster.
  \item The uncompressed size of 30 days of data is 36\TB. What
        fraction of the 30-day working set fits in the LLAP cache?
  \item ORC's columnar layout allows the cache to store only the columns
        accessed by queries rather than full rows. If analyst queries
        reference an average of 8 of the 120 columns in
        \texttt{sales\_events}, recalculate the fraction of the
        working set that fits in cache.
  \item Propose two changes to the LLAP configuration or table design
        that would increase the cache hit rate, and explain the trade-off
        of each.
\end{enumerate}

\bigskip
\exercisesection{Design Problems}

\dpr \textbf{Multi-Engine SQL Platform for a Media Company.}
A digital media company operates a 1{,}000-node Hadoop cluster with
400\TB\ of data. Three teams use the cluster:
\begin{itemize}
  \item \textbf{Data Engineering:} runs nightly ETL pipelines (Hive/MapReduce),
        producing 20 curated tables from 150\TB\ of raw log data.
  \item \textbf{Analytics:} runs 500 ad-hoc queries per day, targeting
        $<$5-minute response time, on the curated tables ($\sim$15\TB).
  \item \textbf{Business Intelligence:} runs 10 recurring dashboard queries
        every 15 minutes on 5 summary tables ($\sim$500\GB total).
\end{itemize}

Design a three-tier SQL platform for this company. Your design must specify:
\begin{enumerate}[label=(\alph*)]
  \item The SQL engine and execution environment (YARN queues, cluster
        sizing) for each team's workload.
  \item How the three engines share the Hive Metastore and Parquet/ORC
        data on HDFS without data duplication.
  \item A query routing service: how incoming queries are classified and
        directed to the correct engine. What metadata does the routing
        service inspect to make this classification?
  \item The LLAP daemon configuration for the BI dashboard tier: how
        many nodes, how much cache, and how to ensure the most-queried
        partitions are warmed in the cache before the first dashboard
        refresh of the day.
  \item A monitoring and SLA framework: which metrics to track for each
        tier, and the alerting thresholds that indicate a tier is degrading.
\end{enumerate}

\dpr \textbf{HiveQL to Presto SQL Migration.}
Your company has 2{,}000 HiveQL queries (stored as \texttt{.hql} files
in a Git repository) that need to be migrated to Presto SQL to reduce
query latency from 30--45 minutes to $<$5 minutes.

Design the migration. Your design must address:
\begin{enumerate}[label=(\alph*)]
  \item The four most common SQL dialect differences between HiveQL and
        Presto SQL that require query rewrites. Give one example of each.
  \item An automated compatibility scanning tool: how to programmatically
        detect which of the 2{,}000 queries require manual rewriting
        vs.\ which can be migrated automatically.
  \item A testing strategy to verify that the migrated Presto queries
        produce results identical to the original HiveQL queries.
        What is the acceptable tolerance for numerical differences?
  \item The migration sequence: which 20\% of the 2{,}000 queries should
        be migrated first, and why? What criterion determines priority?
  \item A rollback plan if a migrated query produces incorrect results
        in production (e.g., a finance report that a business user flags
        as wrong).
\end{enumerate}

\bigskip
\exercisesection{Programming Exercises}

\pe \textbf{Window Function Analytics Pipeline (HiveQL).}
Write a HiveQL script that analyses a table \texttt{daily\_sales}
(columns: \texttt{region STRING}, \texttt{product\_id STRING},
\texttt{sale\_date STRING}, \texttt{revenue DOUBLE}) and produces a
report containing:
(a) 7-day moving average revenue per region-product;
(b) revenue rank within region for each date;
(c) day-over-day revenue growth percentage;
(d) cumulative revenue per region from the start of the month.

\begin{lstlisting}[style=hiveql, caption={PE7.1: Window function analytics pipeline}]
-- Create the source table (run once to set up)
CREATE TABLE IF NOT EXISTS daily_sales (
    region      STRING,
    product_id  STRING,
    sale_date   STRING,
    revenue     DOUBLE
)
STORED AS ORC
PARTITIONED BY (dt STRING);

-- Main analytics query: 4 window metrics in one pass
SELECT
    region,
    product_id,
    sale_date,
    revenue,

    -- (a) 7-day moving average
    ROUND(
        AVG(revenue) OVER (
            PARTITION BY region, product_id
            ORDER BY sale_date
            ROWS BETWEEN 6 PRECEDING AND CURRENT ROW
        ), 2
    )   AS ma7_revenue,

    -- (b) Revenue rank within region for each date
    RANK() OVER (
        PARTITION BY region, sale_date
        ORDER BY revenue DESC
    )   AS daily_rank_in_region,

    -- (c) Day-over-day growth %
    ROUND(
        100.0 * (revenue - LAG(revenue, 1, revenue) OVER (
                     PARTITION BY region, product_id
                     ORDER BY sale_date
                 )) /
        NULLIF(LAG(revenue, 1, revenue) OVER (
                   PARTITION BY region, product_id
                   ORDER BY sale_date
               ), 0),
        2
    )   AS dod_growth_pct,

    -- (d) Month-to-date cumulative revenue per region
    SUM(revenue) OVER (
        PARTITION BY region, SUBSTR(sale_date, 1, 7)
        ORDER BY sale_date
        ROWS BETWEEN UNBOUNDED PRECEDING AND CURRENT ROW
    )   AS mtd_revenue

FROM    daily_sales
WHERE   dt BETWEEN '2024-01-01' AND '2024-03-31'
ORDER BY region, product_id, sale_date;
\end{lstlisting}

\pe \textbf{Presto Query Profiler (Python).}
Implement a Python tool that connects to a Presto/Trino coordinator via
its REST API, submits a SQL query, polls for completion, and reports
execution statistics: elapsed time, bytes scanned, rows processed,
number of splits, and peak memory usage.

\begin{lstlisting}[style=pyspark, caption={PE7.2: Presto query profiler via REST API}]
import time
import requests
import json

PRESTO_HOST = "http://presto-coordinator:8080"
HEADERS     = {
    "X-Presto-User":    "analyst",
    "X-Presto-Catalog": "hive",
    "X-Presto-Schema":  "analytics",
    "Content-Type":     "application/json",
}

def submit_query(sql: str) -> str:
    """Submit a query and return the query ID."""
    resp = requests.post(
        f"{PRESTO_HOST}/v1/statement",
        headers=HEADERS,
        data=sql
    )
    resp.raise_for_status()
    data = resp.json()
    return data.get("nextUri"), data.get("id")

def poll_query(next_uri: str) -> dict:
    """Poll until the query completes; return final stats."""
    while next_uri:
        resp = requests.get(next_uri, headers=HEADERS)
        resp.raise_for_status()
        data = resp.json()
        state = data.get("stats", {}).get("state", "UNKNOWN")

        if state in ("FINISHED", "FAILED", "CANCELED"):
            return data

        # Print progress
        stats = data.get("stats", {})
        pct = stats.get("progressPercentage", 0)
        print(f"  [{state}] {pct:.1f}% complete "
              f"({stats.get('completedSplits',0)}/{stats.get('totalSplits',0)} splits)",
              end="\r")

        next_uri = data.get("nextUri")
        time.sleep(0.5)

    return {}

def profile_query(sql: str):
    """Submit, poll, and report stats for a Presto SQL query."""
    print(f"Submitting: {sql[:80]}{'...' if len(sql)>80 else ''}")
    t_start   = time.time()
    next_uri, query_id = submit_query(sql)
    result    = poll_query(next_uri)
    elapsed   = time.time() - t_start

    stats = result.get("stats", {})
    state = stats.get("state", "UNKNOWN")
    print(f"\n{'='*60}")
    print(f"Query ID     : {query_id}")
    print(f"Final state  : {state}")
    print(f"Elapsed      : {elapsed:.2f}s")

    if state == "FINISHED":
        print(f"Rows read    : {stats.get('processedRows',0):,}")
        print(f"Bytes scanned: {stats.get('processedBytes',0)/1e9:.2f} GB")
        print(f"Splits done  : {stats.get('completedSplits',0):,}")
        print(f"Peak mem     : {stats.get('peakMemoryBytes',0)/1e9:.2f} GB")
        print(f"CPU time     : {stats.get('cpuTimeMillis',0)/1000:.2f}s")
    elif state == "FAILED":
        err = result.get("error", {})
        print(f"Error        : {err.get('errorName','?')}: {err.get('message','?')}")

# --- Run a sample query ---
SQL = """
SELECT  region,
        COUNT(*)          AS num_sales,
        SUM(revenue)      AS total_revenue,
        AVG(revenue)      AS avg_revenue
FROM    hive.analytics.daily_sales
WHERE   dt BETWEEN '2024-01-01' AND '2024-03-31'
GROUP BY region
ORDER BY total_revenue DESC
"""

profile_query(SQL)
\end{lstlisting}

\pe \textbf{HiveQL Complex Type Analytics (HiveQL + Python).}
A table \texttt{session\_events} stores user session data with an
\texttt{ARRAY<STRING>} column \texttt{pages\_visited} (ordered list
of page URLs per session) and a \texttt{MAP<STRING,DOUBLE>} column
\texttt{engagement\_scores} (per-page engagement metrics). Write:
(a) a HiveQL query using \texttt{LATERAL VIEW EXPLODE} to compute
the top-10 most visited page URL patterns; (b) a Python script that
reads the result and visualises the URL frequency distribution.

\begin{lstlisting}[style=hiveql, caption={PE7.3a: LATERAL VIEW EXPLODE on session pages}]
-- Top 10 most visited pages (unnesting the pages_visited array)
SELECT  page_url,
        COUNT(*)                    AS visit_count,
        COUNT(DISTINCT session_id)  AS unique_sessions,
        ROUND(AVG(eng_score), 4)    AS avg_engagement
FROM    session_events
LATERAL VIEW EXPLODE(pages_visited)    pv_table  AS page_url
LATERAL VIEW EXPLODE(engagement_scores) es_table AS eng_page, eng_score
WHERE   dt = '2024-03-01'
  AND   eng_page = page_url          -- correlate map key with array element
GROUP BY page_url
HAVING COUNT(*) > 1000               -- filter low-frequency pages
ORDER BY visit_count DESC
LIMIT 10;
\end{lstlisting}

\begin{lstlisting}[style=pyspark, caption={PE7.3b: Visualise URL frequency from Hive results}]
import pandas as pd
import matplotlib
matplotlib.use("Agg")
import matplotlib.pyplot as plt

# Simulate result of the HiveQL query above
results = {
    "page_url": ["/home", "/search", "/listing/123", "/booking",
                 "/messages", "/profile", "/wishlist", "/help",
                 "/explore", "/trips"],
    "visit_count": [9_812_441, 7_234_190, 4_112_030, 2_981_200,
                     1_822_441, 1_654_022, 1_244_330, 987_441,
                     872_301, 641_200],
    "avg_engagement": [0.72, 0.65, 0.81, 0.89, 0.77, 0.68, 0.74,
                       0.45, 0.58, 0.83],
}
df = pd.DataFrame(results).sort_values("visit_count", ascending=True)

fig, axes = plt.subplots(1, 2, figsize=(14, 6))

# Bar chart: visit counts
axes[0].barh(df["page_url"], df["visit_count"] / 1e6,
             color="#003c78", alpha=0.8)
axes[0].set_xlabel("Visit count (millions)")
axes[0].set_title("Top 10 Pages by Visit Count")
axes[0].grid(axis="x", alpha=0.3)

# Scatter: engagement vs. volume
sc = axes[1].scatter(
    df["visit_count"] / 1e6,
    df["avg_engagement"],
    s=120, c=df["visit_count"] / 1e6,
    cmap="Blues", alpha=0.85, edgecolors="k", linewidths=0.5
)
for _, row in df.iterrows():
    axes[1].annotate(row["page_url"].split("/")[-1] or "home",
                     (row["visit_count"]/1e6, row["avg_engagement"]),
                     textcoords="offset points", xytext=(5, 3),
                     fontsize=8)
axes[1].set_xlabel("Visit count (millions)")
axes[1].set_ylabel("Avg engagement score")
axes[1].set_title("Engagement vs.\ Volume by Page")
axes[1].grid(alpha=0.3)

plt.tight_layout()
plt.savefig("page_analytics.png", dpi=150)
print("Saved: page_analytics.png")
\end{lstlisting}
